% ---------------------------------------------------------------------------
% Evaluation metrics (Methodology, Section 4) - drafted by M1.
% Citation keys match report/references.bib. Numbers come from results/
% (notebooks/02_baselines.ipynb); test-set results are in results_discussion.tex.
% Needs in the preamble: \usepackage{booktabs, amsmath} (table rules, equation).
% ---------------------------------------------------------------------------
\subsection{Evaluation Metrics}
\label{sec:metrics}

The task is closed-set classification of a clip into one of $K=12$ words. Every model outputs a probability
vector $\hat{\mathbf{p}}_i$, and every model is scored by the same code (\texttt{src/evaluate.py}), so the
numbers are comparable across approaches. Table~\ref{tab:metrics} summarises the metrics and what each adds.

\paragraph{Log loss (co-primary)}
The official challenge metric~\cite{zindi2021swahili} is the multi-class log loss,
\begin{equation}
  \mathcal{L} = -\frac{1}{N}\sum_{i=1}^{N} \log \hat{p}_{i,y_i},
\end{equation}
where $y_i$ is the true word of clip $i$. Log loss is a strictly proper scoring rule: its expected value is
minimised only by reporting the true class probabilities, so it rewards honest confidence rather than just the
top-ranked label~\cite{gneiting2007strictly}. This matters for the intended use. A voice interface for mobile
money or a helpline must decide whether to act on a recognised digit or ask the caller to repeat it, and that
decision depends on the confidence. Log loss also punishes confident mistakes severely. Assigning probability
$0.01$ to the true word costs $4.6$ nats, against $\ln 12 \approx 2.48$ for a uniform guess.

Our own experiments show why accuracy alone is not enough. For the MFCC-statistics SVM (an early A1 variant on 40 MFCCs), the probability
method changed the log loss far more than the accuracy. libsvm's pairwise-coupled Platt
probabilities~\cite{platt1999probabilistic,wu2004probability} gave a validation log loss of $0.880$ at $74.3\%$
accuracy. The one-vs-rest sigmoid calibration that scikit-learn now recommends gave $1.230$ at $73.5\%$, and
temperature calibration gave $1.153$ at $73.7\%$. Accuracy is almost unchanged, yet log loss is $40\%$ worse.

\paragraph{Accuracy and macro-F1 (co-primary)}
The training set is exactly balanced (350 clips per word) and the frozen validation and test splits hold 52--53
clips per word. Accuracy is therefore not inflated by a majority class, and it is the conventional headline for
keyword spotting~\cite{warden2018speech}. We also report macro-F1, the unweighted mean of per-class
F1~\cite{sokolova2009systematic}. It penalises a model that over-predicts one word, which accuracy on a balanced
set only partly reveals. We compute it as the mean of per-class F1 scores, not the F1 of mean precision and
recall, since the two definitions can diverge~\cite{opitz2019macro}. Per-class precision, recall and F1, and
the $12\times12$ confusion matrix, show which words fail. For example, A1's recall is only $66\%$ for
\textit{mbili}, which it confuses with \textit{ndio}.

\paragraph{Calibration}
We measure calibration with the expected calibration error (ECE, 15 equal-width confidence
bins)~\cite{naeini2015obtaining,guo2017calibration} and with reliability diagrams. Poor calibration is corrected
post hoc with temperature scaling~\cite{guo2017calibration}: a single scalar $T$ divides the logits (for the
HMM, the per-frame log-likelihoods) before the softmax. It changes the confidence but never the predicted
class, so accuracy and F1 are unaffected. To keep the validation log loss honest, $T$ is \emph{cross-fitted}:
each validation fold is calibrated with a $T$ fitted on the other four folds. $T$ is then refitted on the whole
validation split for use on the test split.

\paragraph{Task-specific diagnostic: sound-alike pairs}
The research question is whether modelling temporal order helps. We therefore track, for word pairs that share
sounds, the share of the pair's clips predicted as the other word of the pair. The pairs are
\textit{tisa}/\textit{sita} (an anagram: the same sounds in a different order), \textit{nne}/\textit{nane},
\textit{tatu}/\textit{tano} and \textit{saba}/\textit{sita}. This links the evaluation directly to the
hypothesis. The order-free A1 confuses \textit{tisa}/\textit{sita} on $6.7\%$ of their validation clips, while
the order-aware HMM (A2) confuses them on none.

\paragraph{Uncertainty and significance}
Model selection uses the validation split only. The test split is used only after every configuration is frozen,
and never for selection. Neural models are trained with three seeds and reported as mean $\pm$ standard
deviation. For every model we compute 95\% percentile bootstrap confidence intervals for accuracy and log loss
(1{,}000 resamples)~\cite{efron1993introduction}; they are listed in \texttt{results/metrics/test\_results.csv}. Pairs of models are compared on the same test clips with the exact
McNemar test on their discordant predictions~\cite{mcnemar1947note}. Dietterich found that this test has an
acceptable Type~I error, and recommends it when each learning algorithm can be run only
once~\cite{dietterich1998approximate}. That fits our setting: retraining the pretrained Transformer the ten
times a $5\times2$ cross-validated test requires is too expensive. The five planned comparisons (the best model
against each other model, and the recurrent A3 against the convolutional A4) are corrected together for
multiplicity with Holm's procedure~\cite{holm1979simple}.

\paragraph{Limitations}
Selection bias, the size of the test split, unknown speakers and the closed vocabulary limit what these metrics can show. They are discussed in Section~\ref{sec:limitations}.

\begin{table}[t]
  \centering
  \caption{Evaluation metrics and their role.}
  \label{tab:metrics}
  \small
  \begin{tabular}{@{}p{0.22\linewidth}p{0.34\linewidth}p{0.34\linewidth}@{}}
    \toprule
    Metric & Why we use it & Limitation \\
    \midrule
    Log loss & Official metric; proper scoring rule; rewards honest confidence & Dominated by a few confident errors \\
    Accuracy & Balanced classes; standard for keyword spotting & Ignores confidence \\
    Macro-F1, per class & Exposes words a model systematically fails & Not a probability metric \\
    ECE + reliability & Directly measures calibration & Depends on binning; sparse bins \\
    Pair confusion & Tests the temporal-order hypothesis & About 105 clips per pair \\
    Bootstrap CI, McNemar & Separates real differences from noise & One split; no speaker grouping \\
    \bottomrule
  \end{tabular}
\end{table}
